\documentclass[11pt, headheight=30pt]{scrartcl}
\usepackage[hw]{darkWhite}

\title{18.755 Lie Algebra}
\subtitle{Homework 1}
\author{Rowechen Zhong with Isaac Zhu}
\begin{document}

\section{Invariant Polynomials}
\begin{prob*}
    Let $R = \CC[x_1, \ldots, x_N, y_1, \ldots, y_N]^{S_N}$ (the algebra of invariant polynomials). Show that $R$ is generated by the elements $Q_{rs} := \sum_{i=1}^N x_i^r y_i^s$ where $1 \leq r + s \leq N$.
\end{prob*}

Let $Q$ be the set generated by $Q_{rs}$.

Consider the elements $z_i = x_i + t y_i$. Then, for all $r \leq N$, $\sum_{i = 1}^N z_i^r = \sum_{i = 1}^N (x_i + t y_i)^r$ is a polynomial in $t$ with coefficients in $Q$ (observe that each term has total degree $r$ in $x$ and $y$). By the usual Newton's identities, we now know that $\sum_{i = 1}^N z_i^r$ is a polynomial $t$ with coefficients in $Q$ for \emph{all} $r$. If we expand this with binomial theorem, we obtain
\[
    z_i^r = \sum_{i = 1}^N \sum_{k = 0}^r \binom{r}{k} x_i^{r-k} t^k y_i^k.    
\]

Now, by roots of unity filter (plug in $t = \omega^j$ for $j = 0,\dots r$ and $\omega$ a $r+1$-th root of unity), we find that each term is linearly independent. (One can also use the Vandermonde determinant trick). Thus, we discover that $\sum_{i = 1}^N x_i^{r-k} y_i^k$ is in $Q$ for all $r,k$. By lemma $27.5$, because $\Delta_N(p)\in Q$ for any $p\in \CC[x,y]$, $Q = R$. Thus we conclude.

\section{Jucys-Murphy Element}
% Let λ = (λ1, . . . , λn) be a partition. Let us fill the Young diagram of λ with numbers, placing c(i, j) := i − j
% in the j-th box in the i-th row. Thus the number written in each box depends only of its position (i, j); it is
% called the content of this box. The content of λ is the sum c(λ) of contents of all its boxes:
% c(λ) = X
% (i,j)∈λ
% c(i, j).
% (i) Show that
% c(λ) = Xn
% i=1
% λi(λi − 2i + 1)/2.
% (ii) Let c =
% P
% 1≤i<j≤N (ij) ∈ C[SN ] be the Jucys-Murphy element (the sum of all transpositions). Show
% that c is a central element of C[SN ] which acts on the irreducible representation πλ of SN by the scalar
% c(λ). (Hint: Consider the action of c on V
% ⊗N and use Schur-Weyl duality to relate it to the diagonal
% action of the quadratic Casimir of gln)
\begin{prob*}
    Let $\lambda = (\lambda_1, \ldots, \lambda_n)$ be a partition. Let us fill the Young diagram of $\lambda$ with numbers, placing $c(i, j) := i - j$ in the $j$-th box in the $i$-th row. Thus the number written in each box depends only of its position $(i, j)$; it is called the content of this box. The content of $\lambda$ is the sum $c(\lambda)$ of contents of all its boxes:
    \[
        c(\lambda) = \sum_{(i,j)\in \lambda} c(i, j).
    \]
    \begin{enumerate}
        \item Show that
        \[
            c(\lambda) = \sum_{i=1}^n \lambda_i(\lambda_i - 2i + 1)/2.
        \]
        \item Let $\mathbf{c} = \sum_{1\leq i < j \leq N} (ij) \in \CC[S_N]$ be the Jucys-Murphy element (the sum of all transpositions). Show that $\mathbf{c}$ is a central element of $\CC[S_N]$ which acts on the irreducible representation $\pi_\lambda$ of $S_N$ by the scalar $(\lambda)$. (Hint: Consider the action of $\mathbf{c}$ on $V^{\otimes N}$ and use Schur-Weyl duality to relate it to the diagonal action of the quadratic Casimir of $\gl_n$)
    \end{enumerate}
\end{prob*}

For part (1), the numbers in the $i$-th row go from $1-i$ to $\lambda_i - i$, which has sum $\lambda_i(\lambda_i - 2i + 1)/2$, so we are done.


For part (2), to show that $\mathbf{c}$ is a central element, it suffices to show that it commutes with some transposition $(pq)$, because transpositions generate $S_N$. Let's break up all the terms in the summation $2\mathbf{c} = \sum_{1\leq i, j \leq N, i\neq j} (ij)$.
\begin{itemize}
    \item If $\{i,j\}\cap \{p,q\} = \emptyset$, then $(ij)(pq) = (pq)(ij)$, so there's nothing to see here.
    \item Suppose $\{i,j\}\cap \{p,q\} = \{p\}$. Then, WLOG $i = p$; the other case is symmetric; and $j = r$. There is another corresponding term, where $i = q$ and $j = r$. These two commute, because
    \[
        (pr)(pq) + (qr)(pq) = (pqr) + (qpr) = (pq)(pr) + (pq)(qr).    
    \]
    \item If $\{i,j\} = \{p,q\}$, then $(pq)(pq) = 1 = (pq)(pq)$.
\end{itemize}
Thus $\mathbf{c}$ commutes with all transpositions, so it is central.

Now, let's use Schur-Weyl duality. We denote the representations of $U(\gl_n)$ and $\CC[S_N]$ on $L_\lambda\otimes \pi_\lambda$ by $A: U(\gl_n)\to \End(L_\lambda \otimes \pi_\lambda)$ and $B: \CC[S_N]\to \End(L_\lambda \otimes \pi_\lambda)$, respectively. We know $B(c)$ is in the centralizer of $\im B$, and thus inside $\im A$; we just want to find a $x\in U(\gl_n)$ such that $A(x) = B(c)$.

It is clear that $B(c)$ is
\[
    \sum_{1\leq i <  j \leq N} \sum_{1\leq k,\ell \leq n} 1 \otimes \cdot \otimes \underbrace{E_{k\ell}}_{i} \otimes \cdot \otimes \underbrace{E_{\ell k}}_{j} \otimes \cdot \otimes 1
\]
by considering its action on pure tensors (this will swap a pair of terms at $i$ and $j$ iff the $i$-th term is $\ell$ and the $j$-th term is $k$, which is what we want). As suggested by the hint, we consider the quadratic Casmir of $\gl_n$ with respect to the trace form $\Tr(xy)$ (a la the vector representation), which is
\[
    \Phi = \sum_{1\leq k,\ell \leq n} E_{k\ell}E_{\ell k}\in U(\gl_n).
\]
\idea{
    [Why one might consider $\Phi$]
    The quadratic Casimir automatically lies in the center of $U(\fg)$, and its action is computable via Lemma 26.6. We would usually default the Killing form, but by the Cartan criterion, we know that the Killing form is degenerate, so it doesn't produce a dual basis or quadratic Casimir. The Killing form is just a special case (adjoint representation) of a more general construction though, which is that this trace form can be associated with any (finite-dimensional) representation.
    
    % A more low-tech solution: stare at $B(c)$ and try to express it as algebraic combinations of $\Delta_n$ things.
}

Then
\begin{align*}
    A(\Omega) &= \sum_{1\leq k,\ell \leq n} A(E_{k\ell})A(E_{\ell k})\\
    &= \sum_{1\leq k,\ell \leq n}\left( \sum_{1\leq i\neq j \leq N} 1\otimes \dots \otimes \underbrace{E_{k\ell}}_i \otimes \dots \otimes \underbrace{E_{\ell k}}_j \otimes \dots \otimes 1\right.\\
    & + \left.\sum_{1\leq i\leq N} 1\otimes \dots \otimes \underbrace{E_{kk}}_i \otimes \dots \otimes 1\right)\\
    &= 2 B(c) + n A(1).
\end{align*}
Thus $B(c) = \frac12\left(A(\Omega) - nA(1)\right)$.

Finally, we have to figure out how $A(\Omega)$ and $A(1)$ act on $L_\lambda$.
First, let's compute the eigenvalue of $\Omega$. Under $\gl_n = \sll_n \oplus \CC I_n$,
\[
    \Omega = \Omega_{\sll_n} + I_n\left(\frac1n I_n\right),
\]
where $\Omega_{\sll_n}$ is the quadratic Casimir of $\sll_n$, as the dual of $I_n$ under the trace form is $\frac1n I_n$. (This is obvious, because the quadratic Casimir is defined as $\sum X_i X^i$ for any basis $\{X_i\}$, and we can simply pick a basis that splits into $\sll_n$ and $I_n$.)

This has eigenvalue $(\lambda_{\sll_n}, \lambda_{\sll_n} + 2\rho_{\sll_n})$ (this is Lemma 26.6). In our case,
\[
    \rho_{\sll_n} = (n-1, n-2, \ldots, 0)
\]
and
\[
    \lambda_{\sll_n} = (\lambda_1, \ldots, \lambda_n)
\]
so
\[
    \lambda_{\sll_n}^* = (\lambda_1 - \frac{|\lambda|}n, \ldots, \lambda_n - \frac{|\lambda|}n)
\].
Therefore, the eigenvalue of $\Omega_{\sll_n}$ is
\begin{align*}
    (\lambda_{\sll_n}, \lambda_{\sll_n}) + 2(\lambda_{\sll_n}, \rho_{\sll_n}) &= \lambda_{\sll_n}(\lambda_{\sll_n}^*) + 2\rho_{\sll_n}(\lambda_{\sll_n}^*)\\
    &= \sum_{i=1}^n \lambda_i^2 - 2i\lambda_i - \frac{|\lambda|^2}n + (n+1)|\lambda|.
\end{align*}
The eigenvalue of $I_n$ is $|\lambda|$, so the eigenvalue of $I_n(\frac1n I_n)$ is $\frac{|\lambda|^2}n$. Hence, the eigenvalue of $\Omega$ is
\[
    \sum_{i=1}^n \lambda_i^2 - 2i\lambda_i + (n+1)|\lambda|.
\]
The eigenvalue of $B(c) = \frac12(A(\Omega) - nA(I_n))$ is
\[
    \frac12\left(\sum_{i=1}^n \lambda_i^2 - 2i\lambda_i + (n+1)|\lambda| - n|\lambda|\right) = \sum_{i=1}^n \frac{\lambda_i(\lambda_i - 2i + 1)}2.
\]
\end{document}